% gz-p8-04-1: 条件概率几何理解——大矩形 Ω + 中间椭圆 A + 上椭圆 B + A∩B
\documentclass[tikz,border=4pt]{standalone}
\usepackage{ctex}
\usepackage{amsmath}
\usetikzlibrary{calc}

\begin{document}
\begin{tikzpicture}[scale=1.0, thick]
  % 大矩形 Ω
  \draw[black, thick] (-4, -2.5) rectangle (4, 2.5);
  \node[font=\small] at (-3.6, 2.2) {$\Omega$};

  % 椭圆 A（中间，较大）
  \fill[red!20] (-0.5, -0.5) ellipse (2.5 and 1.5);

  % 椭圆 B（上方）
  \fill[blue!20] (0.8, 0.8) ellipse (2.2 and 1.3);

  % A ∩ B 交集——用裁剪
  \begin{scope}
    \clip (-0.5, -0.5) ellipse (2.5 and 1.5);
    \fill[purple!50] (0.8, 0.8) ellipse (2.2 and 1.3);
  \end{scope}

  % 重画边界
  \draw[red, thick] (-0.5, -0.5) ellipse (2.5 and 1.5);
  \draw[blue, thick] (0.8, 0.8) ellipse (2.2 and 1.3);

  % 标签
  \node[font=\small, red] at (-2.4, -1.2) {$A$};
  \node[font=\small, blue] at (2.6, 1.6) {$B$};
  \node[font=\footnotesize, purple] at (0.5, 0.0) {$A\cap B$};

  % 注释
  \node[font=\small, below, align=center] at (0, -3.2)
    {$P(B\mid A)=\dfrac{P(A\cap B)}{P(A)}$:};
  \node[font=\footnotesize, below, align=center] at (0, -3.9)
    {在 $A$ 中 $B$ 所占的比例};
\end{tikzpicture}
\end{document}
